\begin{table}[htbp]
\centering
\caption{Salinity Exposure Diagnostics in the Main SFD Yield Samples}
\label{tab:salinity_exposure_diagnostics_supervisor_response}
\small
\resizebox{\textwidth}{!}{%
\begin{tabular}{lrrrrr}
\toprule
Crop & Observations & Municipalities & Pairs & Mean ECe (dS/m) & Above FAO \\
\midrule
Corn & 114,025 & 4,435 & 4,070 & 0.614 & 55.1\% \\
Rice & 66,072 & 3,660 & 3,280 & 0.641 & 49.0\% \\
Cassava & 94,296 & 4,282 & 3,860 & 0.622 & 55.3\% \\
Beans & 102,376 & 4,400 & 4,027 & 0.618 & 56.1\% \\
Soybeans & 33,563 & 2,004 & 1,774 & 0.592 & 7.3\% \\
Sugarcane & 57,069 & 3,536 & 3,117 & 0.623 & 59.9\% \\
\bottomrule
\end{tabular}
}
\par\addvspace{0.5ex}
\parbox{0.95\textwidth}{\textit{Notes:} The table summarizes the estimation samples used in the main crop-yield Spatial First-Difference specifications. Observations are municipality-pair-crop-year observations. Municipalities count unique east- and west-side municipalities observed in each crop sample. Mean ECe is reported in dS/m; the source GeoTIFF rasters are stored as centi-dS/m and the processing divides raw values by 100 before constructing salinity measures. ``Above FAO'' is the share of crop-pair-year observations above the crop-specific FAO agronomic threshold.}
\end{table}
